\chapter{\label{ch:5-conclusion}Conclusion}

In this thesis, we started in \autoref{ch:1-intro} and \ref{ch:2-litreview} 
by exploring the existing approaches to the problem of classical simulation with the ZX-Calculus.
We have seen that the existing stabiliser decomposition algorithms are powerful tools,
as demonstrated in \cite{Kissinger2022cat}.
We also saw that heuristic approaches can be used to improve the performance of these algorithms,
as demonstrated in \cite{sutcliffe2024proccut,Codsi2022}.
However, as was challenged in \cite{sutcliffe2024proccut}, and alluded to in \cite{Codsi2022},
obtaining lower $\alpha$ decompositions need not be just by looking at the immediate $T$-count reduction
from a decomposition, but also by exploiting the structure of the ZX-diagrams 
after carrying out the decomposition.
This involves looking at how the $T$-like spiders can fuse after using the \CutRule decomposition,
based on certain local patterns.

Thus, in \autoref{ch:3-methods} we formalised the heuristic method and applied the formalism 
to the existing approaches.
We also developed new heuristics that take advantage of structure made available in
the \nameref{alg:zx-simplify} algorithm of \cite{Duncan2020graphtheoretic}.
We presented the results of these heuristics in \autoref{ch:4-results} and showed that 
the \texttt{paired+single} heuristic can greatly improve 
the performance of the existing stabiliser decomposition algorithms.

We also explored various classes of circuits and how some of their structures
lend very well to the heuristics developed.
We managed to match the performance of \cite{codsi2023} for IQP circuits,
and showed a clear improvement for circuits with a high number of CCZ gates.
For modified hidden shift circuits, we obtained an improvement of 3 times the $T$-count
that could be simulated within the time limit.

We showed that $\alpha$ is generally reduced for all classes of circuits tested,
though this does not necessarily translate to a reduction in time.
This happens with random Clifford+T circuits coming from Pauli exponentials,
where overhead of computing the heuristics outweighed the saved terms for the range of $T$-counts tested.
However, we have also shown that this overhead is insignificant for circuits with a high $T$-count,
though this remains to be tested for larger circuits and more compute power to match.

We therefore conclude that the heuristic method to stabiliser decomposition
is not only limited to that of the CNOT sandwich \cite{sutcliffe2024proccut},
or the trivial \CutRule decomposition on 2-legged phase gadgets \cite{Codsi2022},
but can also be extended to a wider range of ZX-diagram structures.
In this thesis, this concentrated on the structure created by the reduced gadget form,
as well as the $\PivotRule$ivot rewrite rule,
but potentially more of such structures could be found.


\section{Future Work}

As mentioned in \autoref{sec:other-heuristics}, the \nameref{sec:subgraph-complement} cut
is a potential area for future work, 
perhaps making use of the future graph partioning work in \cite{sutcliffe2024part}.

More compute power could also test the potential of the \nameref{sec:k-connected-heuristics}
on larger circuits, and also to determine the size of the circuit for which the overhead of the heuristics
is outweighed by the saved terms.

Further analysis could also be done on \nameref{sec:complexity} of computing the heuristics,
and whether there could also be a heuristic to not use or only partially use \nameref{alg:zx-simplify}
between each step of a decomposition.
This could allow the overhead of using heuristics to be minimised,
while still providing the $\alpha$-reduction benefits.
Particularly, there are few ways this could be done:
\begin{itemize}
    \item Coming up with a heuristic based version of \nameref{alg:zx-simplify} 
    that knows when \textit{not} to simplify certain vertices to aid the \CutRule decomposition.
    \item Parallelising the heuristic computation for each vertex in the ZX-diagram to avoid 
    the $poly(n)$ time overhead.
    \item Developing much cheaper meta-heuristics that can decide when to use the heuristics.
    \item Bounding how often the heuristics are used or computed in the overall decomposition.
\end{itemize}

Another direction for heuristics is also to also have them for the \catstate{n} decompositions.
These would be more complex due to different resulting patterns in the resulting terms of a decomposition,
but could be very beneficial.
The ideal heuristic would be able to give the best decomposition not just using \CutRule,
but also from any known decomposition.


